\scanpage{56}
\mspage{15}
i.e.
\[
\boxed{r_{ijh}=r_6+\deleted{\unsure{1}}(r_1+\ldots+r_{i-1})+\deleted{\unsure{1}}(r_2\ldots+r_{j-1})
\deleted{\unsure{{}+2r_1+1r_2|2r_1+r_2}}+(r_3\ldots+r_{h-1})}
\]
pour $i<j<h$
\[
\left(\begin{aligned}\text{poids }&i+j+h-5\\\text{maximum }&10\end{aligned}\right)
\quad
\deleted{\begin{gathered}\text{poids }1+3(i-1)+2(j-i\unsure{{}+1|{}-1})+h-j\unsure{{}+1}\\
=1+(i+j+h\unsure{{}-5})=i+j+h-5\\\text{max. }10\end{gathered}}
\]
enfin
\[
\begin{aligned}
r_0&=2\eta-\textstyle\sum\zeta_i=(\eta-\zeta_1-\zeta_2-\zeta_3)+(\eta-\zeta_4-\zeta_5-\zeta_6)\\
\text{i.e.}\quad &=\underbrace{r_{123}}_{\deleted{\unsure{r_0+2r_1+r_2|r_6+2r_1+r_2}}}+r_{456}
\end{aligned}
\]
donc
\[
\boxed{r_0=\deleted{\unsure{r_6+2r_1}}\ r_1+2r_2+3r_3+2r_4+r_5+2r_6}
\qquad(\text{poids }\deleted{\unsure{14}}11)
\]
(NB On a là la ``plus grande racine'' \deleted{\unsure{évidt}}).

Donc tte racine est \deleted{somme} comb. lin. à coeff. entiers tous de $=$ signe de $r_1\ldots r_6$, qui forment donc un syst. de racines simples. Enfin, on vérifie immédiatement que $r_i.r_j=0$ \inserted{pour $i<j$} sauf dans les cas \deleted{$j\ne6$} $j=i+1\deleted{\unsure{{}\ne}} \le5$ et \unsure{le p.|les p.|le cas} $j=6,\ i=3$, ce qui montre que le diagramme est $E_6$. ---

\uline{Corollaire} Soit $\widetilde W$ le groupe des autom du syst. de racines $(E(\Delta),R(\Delta))$, et considérons l'application canonique
\[
W=\operatorname{Aut}(\Delta)\lhook\joinrel\longrightarrow\widetilde W
\]
Alors \unsure{la} l'image est d'indice $2$ dans $\widetilde W$.

En effet, on sait que $\widetilde W/W\simeq$ groupe des autom. du diagramme de Dynkin, qui est ici d'ordre $2$.
% archive sheet number: 55
% archive sheet number: 56 (at the far right of the facing-page edge)
